\subsection{Descripción del Sistema}
La app \textbf{Packlead} es una aplicación móvil pensada para seguir los pedidos en tiempo real, sobre todo para mejorar todo lo que tiene que ver con la entrega de última milla en logística.
Está hecha con una arquitectura modular, o sea que separa bien la lógica del negocio (las reglas y procesos importantes) de la parte técnica de infraestructura. Así se maneja fácil los diferentes roles: por un lado los administradores y por otro los repartidores o dispatchers. Esta forma de organizarlo hace que el sistema sea más fácil de escalar, de mantener a largo plazo y de entender quién hace qué en la plataforma.

\subsubsection{Propósito}
El propósito principal de Packlead es ser una solución completa para manejar la logística de entregas, haciendo más fácil controlar todo lo que pasa en las operaciones y tener siempre una vista clara de los procesos. La plataforma te deja gestionar desde un solo lugar tanto los pedidos como a los repartidores, usando una interfaz pensada especialmente para el control.

Además, ofrece visibilidad en tiempo real del estado de los pedidos y de la ubicación de los repartidores, ayudando a optimizar los tiempos de entrega y a tomar decisiones rápidas en el día a día. Por último, el sistema mantiene todo sincronizado entre el servidor y los dispositivos móviles gracias a tecnologías en la nube, garantizando que la información esté siempre actualizada y consistente para todos.

\subsubsection{Alcance}
La aplicación incluye varias funcionalidades clave para manejar de forma completa los pedidos y la distribución. Entre las más importantes están:

\begin{itemize}
    \item El \textbf{Módulo administrativo}, que permite crear y gestionar a los repartidores, generar órdenes nuevas y ver indicadores clave de desempeño con gráficos fáciles de entender.
    \item El \textbf{Seguimiento geográfico}, que integra mapas interactivos para mostrar las rutas y los puntos de entrega en tiempo real.
    \item La \textbf{Gestión del ciclo de vida} de las órdenes, que controla con detalle los diferentes estados por los que pasa cada pedido, como pendiente, en camino o entregado.
    \item El \textbf{Esquema multiperfil}, que maneja la autenticación de forma diferente según el rol del usuario y adapta los flujos de trabajo a cada tipo de perfil.
\end{itemize}

\subsubsection{Límites y Exclusiones}
En su configuración actual, el sistema tiene algunas limitaciones importantes que hay que tener en cuenta al usarlo. Todo depende mucho de servicios de terceros, sobre todo de Google Maps Platform y Firebase, así que si alguno de estos falla o se cae, la app puede dejar de funcionar correctamente o interrumpirse por completo.
Además, como usa Firebase Realtime Database, necesita que el celular tenga internet todo el tiempo para actualizar las coordenadas y sincronizar la información entre el servidor y los repartidores. El seguimiento de la ubicación se basa solo en el GPS del teléfono del repartidor, por lo que la precisión puede variar dependiendo del modelo del celular o de las condiciones del momento. Por último, en la versión actual no hay integración de pasarelas de pago ni nada para manejar transacciones financieras directamente desde la plataforma.

\subsubsection{Componentes Principales}
El sistema está organizado en cuatro componentes principales que se ven bien claros en cómo está hecho técnicamente.

\begin{itemize}
    \item El primero es el \textbf{núcleo de gestión de estado}, que usa Riverpod para manejar el flujo de los datos y para que la interfaz reaccione rápido cada vez que algo cambia.
    \item El segundo es la \textbf{capa de servicios de ubicación}, que combina Geolocator para sacar las coordenadas del GPS y Google Maps para mostrar todo en un mapa de forma visual.
    \item El tercero es el \textbf{motor de comunicación}, que se apoya en librerías como Dio y Http para hacer las peticiones asíncronas a las APIs externas sin que la app se trabe.
    \item El cuarto es la \textbf{infraestructura en la nube}, basada en Firebase, que sirve como el lugar central donde se guardan las rutas en tiempo real y se sincronizan todos los datos entre el servidor y los móviles.
\end{itemize}

\subsection{Casos de Uso}
Para dejar bien claro cómo interactúan los usuarios con el sistema Packlead, se definen varios casos de uso. Estos casos explican las funciones principales que tienen los administradores y los repartidores (o dispatchers), enfocándose en las operaciones más importantes para manejar la logística, seguir los pedidos y controlar todo el proceso de entregas.

Definir estos casos de uso ayuda a organizar los requisitos funcionales del sistema, facilita analizarlo mejor y sirve de guía sólida para el diseño, la programación y las pruebas que se hagan después.

\begin{table}[H]
\centering
\caption{Matriz de Casos de Uso de Packlead}
\vspace{3mm}
\begin{tabularx}{\textwidth}{l l L}
\toprule
\textbf{Código} & \textbf{Nombre del Caso de Uso} & \textbf{Descripción Técnica} \\
\midrule
UC-01 & Autenticar Usuario & Permite a administradores y dispatchers acceder al sistema mediante validación de credenciales utilizando el módulo \texttt{auth\_provider.dart}. \\
\addlinespace
UC-02 & Registrar Pedido & El administrador ingresa los datos del pedido (cliente, teléfono, zona, fecha y ubicación), inicializándolo en estado \textit{pending}. \\
\addlinespace
UC-03 & Asignar Dispatcher & El administrador asocia un \texttt{dispatcherId} a un pedido existente para habilitar su visualización en la hoja de ruta. \\
\addlinespace
UC-04 & Gestionar Repartidores & Permite crear, editar y visualizar perfiles, así como monitorear su estado (\textit{available/inactive}). \\
\addlinespace
UC-05 & Consultar Hoja de Ruta & El dispatcher visualiza los pedidos asignados pendientes de entrega según su identificador. \\
\addlinespace
UC-06 & Actualizar Estado & El dispatcher actualiza el estado del pedido siguiendo el flujo \textit{pending $\rightarrow$ shipped $\rightarrow$ delivered}. \\
\addlinespace
UC-07 & Cambiar Disponibilidad & El dispatcher puede modificar su estado operativo para recibir o pausar asignaciones. \\
\addlinespace
UC-08 & Visualizar Seguimiento & El administrador supervisa en tiempo real o mediante historial la ubicación y progreso de las entregas. \\
\bottomrule
\end{tabularx}
\end{table}

\subsection{Reglas de Negocio}
Las reglas de negocio establecen restricciones operativas que garantizan la coherencia del proceso logístico, la integridad de los datos y la correcta aplicación de los roles dentro del sistema.

\begin{table}[H]
\centering
\caption{Reglas de Negocio del Sistema Logístico}
\vspace{3mm}
\begin{tabularx}{\textwidth}{l p{4cm} L}
\toprule
\textbf{Código} & \textbf{Regla de Negocio} & \textbf{Descripción} \\
\midrule
RN-01 & Ciclo de Vida del Pedido & Cada pedido solo puede adoptar tres estados: pendiente, en ruta y entregado, definidos mediante el \texttt{enum OrderState}. \\
RN-02 & Asignación de Roles & Segregación estricta según \texttt{UserRole} (admin, dispatcher, none). Rol \textit{none} no posee acceso operativo. \\
RN-03 & Geolocalización Obligatoria & No se permite registrar pedidos sin coordenadas válidas (latitud/longitud). \\
RN-04 & Estado de Repartidor & Repartidores en estado inactivo no son elegibles para nuevas asignaciones. \\
RN-05 & Integridad Temporal & Almacenamiento de fechas en formato ISO 8601 para interoperabilidad entre Flutter y Backend. \\
\bottomrule
\end{tabularx}
\end{table}

\subsection{Requisitos No Funcionales}
Los requisitos no funcionales definen atributos de calidad y restricciones técnicas que influyen directamente en el rendimiento, seguridad y mantenibilidad.

\begin{table}[H]
\centering
\caption{Especificación de Requisitos No Funcionales}
\vspace{3mm}
\begin{tabularx}{\textwidth}{l l L}
\toprule
\textbf{Código} & \textbf{Categoría} & \textbf{Requisito} \\
\midrule
RNF-01 & Arquitectura & Implementación reactiva basada en \textbf{Riverpod} para gestión de estado desacoplada. \\
RNF-02 & Portabilidad & Ejecución multiplataforma en Android e iOS mediante el framework \textbf{Flutter}. \\
RNF-03 & Persistencia & Integración con \textbf{Firebase} para sincronización en tiempo real y almacenamiento \textit{cloud}. \\
RNF-04 & Usabilidad / UI & Soporte dinámico para modo claro y oscuro según la configuración del SO. \\
RNF-05 & Configuración Segura & Uso de archivos \texttt{.env} y scripts para manejo de claves API, evitando \textit{hardcoding}. \\
RNF-06 & Robustez & Manejo de rutas inexistentes mediante pantalla \texttt{ScreenNotFound} para evitar cierres inesperados. \\
\bottomrule
\end{tabularx}
\end{table}